\scanpage{376}
\mspage{115}
% archive sheet number: 375

Dans le cas géométrique $T=\unsure{\underline{\mathrm{Tg}}|\underline{T_g}}$, et \unsure{d'} $C$ de la forme $\underline{\mathrm{Ch}}_{\unsure{\Sigma|\Sigma_{\cdot}}}(X;a,b)$, donnant lieu à $C_\infty(a,b)$, on \uline{n'a pas} de définition naturelle d'un \inserted{\unsure{bi-|bi}}foncteur de composition des chemins
\[
C_\infty(a,b)\times C_\infty(b,c)\longrightarrow C_\infty(a,c).
\]

On peut définir, certes, des accouplements
\[
\underline{\mathrm{Ch}}_{\Sigma,\unsure{\tau|z}}(a,b)\times
\underline{\mathrm{Ch}}_{\Sigma,\tau'}(b,c)\mathrel{\unsure{-|\to}}
\underline{\mathrm{Ch}}_{\Sigma,\tau\tau'}(a,c),
\]
\deleted{mais} ou, avec les notations relatives à la donnée de
\[
\underline{t}:\mathbb N^\circ\mathrel{-}\deleted{\unsure{C}}\underline{\mathrm{Tg}}
\]
\[
\underline{\mathrm{Ch}}_n(a,b)\times\underline{\mathrm{Ch}}_m(b,c)\mathrel{-}
\underline{\mathrm{Ch}}_{\unsure{n+m|n*m}}(a,c),
\]
mais ces accouplements \uline{ne sont pas} compatibles avec les morphismes de transition des trois syst. inductifs de catégories $(\mathrm{Ch}_n)_{n\in\mathbb N}$, ils ne le sont \uline{qu'à homotopie près}. Cela suffit cependant pour avoir des \unsure{accouplement|accouplements} parfaits
\[
H_T(a,b)\times H_{\unsure{T|t}}(b,c)\mathrel{-}H_T(a,c),
\]
avec associativité, et \inserted{(si $\Sigma=\mathrm{Fl}(X)$ tout entier)} l'existence de \uline{vrais} inverses, i.e. de \deleted{foncteurs} flèches (dans $\mathrm{Ind}\,\deleted{\unsure{\mathrm{Hot}}}\mathrm{Cat\,hot}$) de ``passage au chemin opposé''
